In order to verify convergence of the series in \eqref{eq:suffAH1}, we shall estimate the sequences $(c_n)$, defined in~\eqref{eq:c_k}, and $(s_{2n})$ in terms of the quantities
\begin{gather}\label{eq:Un}
U_n:=\frac{s_{2n}}{b_0^2b_1^2\cdots b_{n-1}^2},\qquad n\ge 1,\qquad U_0:=s_0=1
\end{gather}
and
\begin{gather}\label{eq:Vn}
V_n:=b_0 b_1\cdots b_{n-1}c_n,\qquad n\ge 1,\qquad V_0:=c_0>1.
\end{gather}
This will be done in the next two sections.

The quantity $U_n$ is defined for any moment problem, while $V_n$ only makes sense for indeterminate problems.

It is easy to see that both quantities are scale invariant in the sense that if $(a_n)$, $(b_n)$ from the recurrence relation \eqref{eq:rec} are replaced by a constant multiple $(\lambda a_n)$, $(\lambda b_n)$ for some $\lambda>0$, then $U_n$, $V_n$ remain unchanged.

In the following we shall often compare quantities like $s_n,U_n,c_n,\ldots$ depending on the recurrence coefficients $(a_n)$, $(b_n)$ with those corresponding to other recurrence coefficients, say $(\tilde{a}_n)$,~$\big(\tilde{b}_n\big)$. We then denote the corresponding quantities $\tilde{s}_n,\tilde{U}_n,\tilde{c}_n,\ldots$.

\begin{lem}\label{thm:Un}\quad
\begin{enumerate}\itemsep=0pt
\item[$(i)$] For any moment problem we have $U_n\ge 1$.
 \item[$(ii)$] For any indeterminate moment problem we have $V_n\ge 1$ and $V_n\sqrt{U_n}=c_n\sqrt{s_{2n}}\ge 1$.
\end{enumerate}
\end{lem}

\begin{proof}By \eqref{eq:x} we have
\begin{gather*}
s_{2n}=\langle x^n,x^n\rangle=c_{0,n}^2+c_{1,n}^2+\dots +c_{n,n}^2 \ge c_{n,n}^2= b_0^2b_1^2\cdots b_{n-1}^2,
\end{gather*}
hence by \eqref{eq:c_k} and \eqref{eq:b_n}
\begin{gather*}
c_n=\left (\sum_{j=n}^\infty b_{n,j}^2\right )^{1/2}\ge b_{n,n}={1\over b_0b_1\cdots b_{n-1}} \ge 1/\sqrt{s_{2n}}.\tag*{\qed}
\end{gather*}\renewcommand{\qed}{}
\end{proof}

In \cite[Section 3]{B:S2} we considered the power series
\begin{gather*}
\Phi(z)=\sum_{n=0}^\infty c_nz^n,
\end{gather*}
and proved that $\Phi$ is an entire function of minimal exponential type. In terms of $(c_n)$ this
can be expressed
\begin{gather}\label{eq:minexpt}
\lim_{k\to\infty} k c_k^{1/k} =0.
\end{gather}
It was also proved there that $\Phi$ has the same order $\rho$ and type $\tau$ as the moment problem, which by definition is the common order and type of the entire functions in the Nevanlinna matrix, see \cite{B:P,B:P:H}. From the general formula for type of an entire function in terms of its power series coefficients, see~\cite{Le}, we then get the following result.
